\documentclass[11pt]{article}
\usepackage{setspace}
\usepackage{geometry}
\geometry{margin=1in}
\setstretch{1.0}

\begin{document}

\noindent
\textbf{Zhongli Wang} \\
Phone: +49 160 9541 6015 \\
Email: nkwangzhongli@gmail.com \\
Website: https://sites.google.com/view/zhongli-wang \\
Bielefeld, Germany \\

\vspace{0.3cm}

\noindent
\textbf{Prof.\ Dr.\ Julian Hinz} \\
Kiel Institute for the World Economy \\
Chausseestraße 111, 10115 Berlin \\

\vspace{0.3cm}

\noindent
\textbf{Application for PostDoc Position in Quantitative Economics (m/f/d) --- AI and Structural Economic Models for Economic Policy}\\

\vspace{0.2cm}

\noindent Dear Professor Hinz,\\[0.15cm]

I am writing to apply for the PostDoc position in Quantitative Economics at the Kiel Institute's Berlin office. I completed my Ph.D.\ in Economics at Bielefeld University and Universit\`a Cattolica del Sacro Cuore (Marie Sk{\l}odowska--Curie EPOC network) in September 2026 and will therefore hold my doctorate before the contract starts in autumn 2027. My research builds structural, simulation-based models of markets with heterogeneous agents and strategic interaction, and estimates, calibrates, and validates them against empirical data.

My dissertation, \emph{Strategic Behavior and Economic Dynamics in ICT Markets: An Agent-Based Framework}, develops models in which firms and heterogeneous users interact under network effects, congestion constraints, and imperfect information, and uses them for counterfactual analysis of regulation in digital markets. Its empirical chapter is devoted to parameter estimation and model validation: it compares the Simulated Method of Moments, Bayesian Monte Carlo, and neural-network surrogate methods, calibrating a model of the ICT industry to financial data of AT\&T and T-Mobile. I program primarily in Python and Julia (also R and MATLAB) and follow good software practice: Git version control, automated testing, and reproducible pipelines.

The project's focus on assessing how reliable AI-assisted model analysis is connects directly with my recent work. In \emph{Agentic Workflows for Economic Research} (with Dawid, Harting, and co-authors, 2025), we designed a framework integrating large-language-model agents into empirical research pipelines (document collection and classification), with explicit quality controls for when LLM output can be trusted. In my daily research, I use Claude Code to build reproducible, well-tested research code. In \emph{Polytope Fraud Theory} (\emph{International Review of Financial Analysis}, 2025), we combined economic reasoning with supervised machine learning (XGBoost) and extensive robustness checks to detect financial-statement fraud. A recurring theme is the question at the heart of the transfer project: when AI-assisted analyses deliver reliable, economically interpretable results.

I am also eager to work at the boundary of research and policy. My doctoral training in the EPOC network (\emph{Economic Policy in Complex Environments}) emphasized policy-relevant modeling, and internships at the People's Bank of China exposed me to policy-institution practice; I also speak German (C1). Beyond the project, I would pursue my own agenda on the economics of digital markets, regulation, and AI-assisted methods, targeting publications in leading international journals.

Thank you for considering my application. I would welcome the opportunity to discuss my fit for the project.

\vspace{0.15cm}

\noindent Sincerely,\\[0.4cm]
\noindent Zhongli Wang

\end{document}